\documentclass[runningheads]{llncs}
\usepackage{lmodern}
\usepackage[T1]{fontenc}
\usepackage{microtype}
\usepackage{amsmath,amssymb}
\usepackage{booktabs}
\usepackage{url}
\usepackage{listings}
\usepackage{xcolor}

\lstset{
  basicstyle=\ttfamily\scriptsize,
  breaklines=true,
  frame=single,
  columns=fullflexible,
  captionpos=b,
  backgroundcolor=\color{gray!5},
  rulecolor=\color{gray!40},
  xleftmargin=3pt,
  xrightmargin=3pt
}

\begin{document}

\title{Supplementary Material: Prompt Templates for QEG-RAG Two-Tier Diagnostic Engine}
\titlerunning{Prompt Templates for QEG-RAG}
\author{Trung-Kien Nguyen\inst{1} \and
Khac-Hoai-Nam Bui\inst{2} \and
Hoang-Hiep Lai\inst{2} \and
Quang-Huy Phung\inst{2} \and
Van-Vinh Nguyen\inst{1}\thanks{Corresponding author.}}
\authorrunning{T.-K. Nguyen et al.}
\institute{VNU University of Engineering and Technology, Vietnam National University, Hanoi, Vietnam\\ \and
Viettel AI, Viettel Group, Hanoi, Vietnam}
\maketitle

\appendix
\section{Prompt Templates for Two-Tier Reasoning Engine}
\label{sec:prompt_templates}

This supplementary document provides the exact prompt templates used in the Two-Tier Controlled Reasoning Engine of QEG-RAG.

\subsection{Tier 1: Primary Reader Prompt Template}
The Tier 1 prompt serializes the compact evidence graph $\mathcal{X}^*(G_q)$, incorporating star record blocks, deduplicated context passages, and deterministic path scaffolds. It instructs the language model to perform immediate fact extraction or trigger a diagnostic handoff.

\begin{lstlisting}[caption={Tier 1 Prompt Template (Primary Reader with Diagnostic Protocol)}]
Answer the question using the grounded compact OUEG payload.

`rows` retains table rows selected by grounded paths, their original order, values, passage references, and edge IDs. `edges` uses the legend schema and contains the other selected KG edges. `contexts` are deduplicated by passage ID. `paths` are reading guides only. Their `score` is rounded relevance metadata, not factual proof. Follow `ordering` when present; use `edges` and `contexts` for multi-hop evidence. Copy the shortest answer span exactly as written in the evidence; do not translate or expand names. Return only that span.

Fallback protocol (only when you cannot prove one answer from this payload): keep the normal successful-answer behaviour above unchanged---return its shortest answer span and nothing else. If you cannot answer, never return a blank response. Instead return exactly one of:
1. `NEED_FOCUS: {"candidates":["candidate names already in the payload"],"blocker_type":"missing_link","bottleneck":"why the answer is not yet unique","needed_check":"exact row, relation, year, rank, or passage fact to compare"}` when a focused L2 re-read can resolve it.
2. `UNRESOLVED: {"blocker_type":"insufficient_evidence","bottleneck":"what evidence is absent or contradictory","needed_check":"what fact would be required"}` when this payload cannot prove an answer.

`blocker_type` must be exactly one label: `candidate_collision`, `missing_link`, `insufficient_evidence`, `conflicting_evidence`, or `answer_type_uncertain`; never combine labels with `|`.
Do not invent candidates or evidence.

Question: {question}
Compact OUEG payload: {compact_json}
Answer:
\end{lstlisting}

\subsection{Tier 2: Targeted Reader Prompt Template}
When Tier 1 returns a \texttt{NEED\_FOCUS} diagnostic signal, the targeted bridge retriever augments missing passages for disputed candidates without calling the LLM. Tier 2 then receives a minimal state capsule $\mathcal{K}_2$ containing \emph{only} the disputed candidates and new bridge evidence, filtering out all unrelated table rows.

\begin{lstlisting}[caption={Tier 2 Prompt Template (Targeted Reader over Minimal State Capsule)}]
Resolve the question using only the grounded L2 state capsule below. This is a new request: do not assume you remember L1 or use facts absent from this capsule.

The capsule contains the original question, the unresolved constraint, complete row bindings for the focus candidates, and their relevant evidence. `coverage` says whether retrieval found new text; absence of evidence is not a negative fact. Scores are intentionally omitted.

Compare the candidates against the evidence to resolve the question. Return:
`FINAL_ANSWER: <shortest evidence span>`
`EVIDENCE_PATH: {"edges":["edge IDs in this capsule"],"passages":["source IDs in this capsule"]}`
for the proven candidate. 

Only return `UNRESOLVED: {"blocker_type":"...","bottleneck":"...","needed_check":"..."}` if the evidence strictly contradicts all candidates. Never return blank or invent evidence.

Question: {question}
L2 state capsule: {compact_json}
Answer:
\end{lstlisting}

\end{document}
